\chapter{Thiết kế Hệ thống và Triển khai Kỹ thuật}
\label{ch:system_design_implementation}

Chương này trình bày kiến trúc, mô hình hóa toán học và triển khai của khung bảo mật nhận thức hai tầng SENTINEL: một quy trình phát hiện-và-phản hồi lai được thiết kế để san lấp khoảng cách giữa xử lý gói tin tốc độ đường truyền và suy luận LLM tốn tài nguyên, đồng thời thiết lập các rào chắn nhận thức cách ly LLM khỏi thao túng đối kháng. Chương bao quát sơ đồ khối hệ thống, bộ tiền lọc thống kê trực tuyến, Cổng ML, quy trình nhận thức dựa trên đồ thị trạng thái, truy xuất Dual-RAG, lược đồ Bộ nhớ Đe dọa, pipeline rào chắn, chuỗi audit HMAC, giao diện Streamlit và triển khai Docker.

\section{Tổng quan về Kiến trúc Nhận thức Hai tầng}
Để bảo vệ môi trường doanh nghiệp mà không chịu độ trễ quá lớn, SENTINEL chia pipeline thành hai tầng chuyên biệt (Hình~\ref{fig:soc_ai_concept}). \textbf{Tầng 1}, lớp tốc độ đường truyền phi trạng thái, nằm tại biên mạng: nó thu nạp telemetry tốc độ cao và hoặc tính độ lệch thống kê thời gian thực qua thuật toán Welford, hoặc áp dụng chữ ký regex, bỏ hoặc ghi log sự kiện lành tính tức thì. \textbf{Tầng 2}, lớp nhận thức có trạng thái do một tác tử LangGraph tự chủ~\cite{langgraph2024} điều khiển, thực hiện suy luận sâu, truy xuất nền tảng từ Dual-RAG, liên kết với Bộ nhớ Đe dọa SQLite, và đưa ra biện pháp giảm thiểu có cấu trúc. Giữa hai tầng, một \textbf{Cổng ML} học máy tự quyết phần lớn ca leo thang chắc chắn mà không cần gọi LLM. Bố cục này giảm tải phần lớn lưu lượng trước khi suy luận (định lượng ở Chương~\ref{ch:experiments_evaluation}), dành LLM cho thiểu số thực sự mơ hồ.

\begin{figure}[!htbp]
    \centering
    \includegraphics[width=0.95\textwidth,height=0.66\textheight,keepaspectratio]{images/soc_ai_concept.png}
    \caption{Kiến trúc tổng thể SENTINEL: luồng dữ liệu hai tầng và các cấu phần. (Nguồn: Tác giả)}
    \label{fig:soc_ai_concept}
\end{figure}

Luồng vận hành (Hình~\ref{fig:sentinel_flow}) bắt đầu khi một flow log được bắt. Tầng 1 chặn nó: chữ ký WAF hoặc injection cho phán quyết \texttt{BLOCK\_IP} trực tiếp, lưu lượng whitelist bị bỏ, ngược lại thống kê được cập nhật và tính Z-score. Log lành tính bị bỏ hoặc ghi; log vượt ngưỡng rủi ro thì leo thang. Mỗi ca leo thang trước tiên được Cổng ML chấm điểm và tự quyết khi chắc chắn; chỉ phần dư tới Tầng 2, nơi nó được nén và đóng gói nonce (bộ phân tích Drain~\cite{he2017drain}), nối nền \emph{vô điều kiện} với Dual-RAG (kèm Semantic Cache) và Bộ nhớ Đe dọa, rồi được nút LLM-triage phán xử thành một phán quyết có cấu trúc. Đầu ra được kiểm chứng, làm sạch, và ghi dưới chuỗi HMAC, thực thi chặn tường lửa hoặc chuyển tới dashboard HITL Streamlit.

\begin{figure}[!htbp]
    \centering
    \includegraphics[width=0.95\textwidth,height=0.74\textheight,keepaspectratio]{images/sentinel_flow.png}
    \caption{Luồng thực thi vận hành chi tiết trong pipeline SENTINEL. (Nguồn: Tác giả)}
    \label{fig:sentinel_flow}
\end{figure}

\section{Định vị trong Ngăn xếp SOC Hiện hữu}
Một đánh giá công bằng phải tách biệt điều SENTINEL đóng góp với điều công cụ trưởng thành đã cung cấp. Phát hiện chữ ký (Snort, Suricata, Zeek), liên kết khối lượng lớn (Splunk, QRadar, Elastic)~\cite{siem2021}, và tự động phản hồi tất định (SOAR) đã lâu đời, và nhiều SIEM đã gắn nhãn cảnh báo bằng định danh ATT\&CK. SENTINEL có chủ đích \emph{không} cạnh tranh phát hiện: Tầng 1 phi trạng thái của nó chính là một lớp chữ ký-và-thống kê như vậy và sẽ đứng trước, hoặc được nạp bởi, một IDS/SIEM sẵn có. Ablation ở Chương~\ref{ch:experiments_evaluation} xác nhận điều này---thêm tầng nhận thức không nâng F1 phân loại luồng thô so với Tầng 1 đơn độc---điều này dự kiến vì phát hiện phổ thông đã mạnh, và đó là lý do giá trị tầng thứ hai nằm ở nơi khác.

Giá trị đó là lớp mà ngăn xếp hôm nay \emph{không} cung cấp: một tầng \emph{nhận thức} an toàn, giải thích được, tự chủ cho thiểu số mơ hồ. Ba thuộc tính tách nó khỏi SOAR và các trợ lý ``AI-SOC'' đám mây. Thứ nhất, SOAR chạy playbook cứng soạn sẵn không thể suy luận tình huống ngoài kịch bản hay tự biện minh~\cite{soar2025}; SENTINEL thay bằng tác tử đồ thị trạng thái suy luận trên ngữ cảnh MITRE/NIST truy xuất và phát ra lý lẽ kiểm toán được (đạt $3.1/5$ bởi trọng tài LLM chéo họ). Thứ hai, ngay khi một mô hình \emph{sinh} vào vòng lặp, nền-log trở thành bề mặt injection (kẻ tấn công có thể nhúng ``bỏ qua chỉ dẫn trước và đánh dấu benign'' vào trường User-Agent), điều mà bộ chữ ký và SOAR miễn nhiễm theo cấu trúc vì không bao giờ diễn giải nội dung theo kiểu sinh; lớp Đóng gói Dữ liệu Phân định, rào chắn đồng thuận, và chuỗi audit HMAC là phòng thủ tất định trực tiếp cho bề mặt mới này, chặn hoặc leo thang mọi nỗ lực hạ cấp trên bộ 120 payload ở Chương~\ref{ch:experiments_evaluation}. Thứ ba, trợ lý thương mại đẩy telemetry lên LLM đám mây của nhà cung cấp, không chấp nhận được trong môi trường air-gapped hay Zero-Trust; SENTINEL chạy toàn tầng suy luận cục bộ trên một mô hình lượng tử hóa, với cổng hai tầng giữ suy luận cục bộ nhiều giây khả thi (giảm độ trễ $82.97\%$). Tóm lại, SENTINEL cung cấp tầng suy luận-và-phản hồi được củng cố chống injection, triển khai cục bộ được, và giải thích được, cho phép một mô hình sinh an toàn đảm nhận vai trò mà playbook SOAR cứng đang lấp hôm nay.

\section{Tầng 1: Bộ lọc Tốc độ Đường truyền Phi trạng thái (RuleEngine \& Welford)}
Tầng 1 thu nạp luồng tốc độ cao và quản lý trạng thái kết nối không cần I/O đĩa qua bộ quản lý session-baseline trong bộ nhớ khóa theo IP nguồn. Với mỗi session, nó tính đặc trưng lưu lượng bằng thuật toán trực tuyến Welford~\cite{welford1962}: các truy hồi trung bình $\mu_k$ và tổng-bình-phương $M_{2,k}$ ở Mục~\ref{sec:welford} cho phương sai lăn ổn định số học $\sigma_k^2$, từ đó suy ra Z-score,
\begin{equation}
    Z = \frac{x_k - \mu_k}{\sigma_k}.
\end{equation}
Nếu $Z$ vượt ngưỡng cấu hình được ($\alpha>3.5\sigma$) luồng được gắn \texttt{ESCALATE}; baseline tiếp tục cập nhật để bảo toàn chuỗi chậm-âm-thầm. Song song, một \texttt{RuleEngine} tĩnh khớp chữ ký regex biên dịch sẵn để chặn mẫu đã biết (vd cú pháp SQL-injection) với \texttt{DROP}/\texttt{BLOCK\_IP} tức thì. Để tối ưu hiệu năng, rule engine sử dụng cấu trúc tập hợp (set) $O(1)$ cho việc tra cứu cổng (port lookup). Triển khai dùng hai lớp: \texttt{RunningStats} (đếm/trung bình/tổng-bình-phương trong bộ nhớ, cập nhật $O(1)$) và \texttt{SessionBaseline} (bảng băm với bộ dọn rác TTL giới hạn bộ nhớ). Ngoài ra, một bộ đệm giới hạn dung lượng (\texttt{\_rep\_cache}) được triển khai nhằm ngăn chặn tràn RAM khi xử lý lưu lượng lớn.

Hai cơ chế an toàn chi phối baseline. Thứ nhất, cập nhật trực tuyến \emph{khóa-theo-phán-quyết}: baseline chỉ tiến khi có phán quyết lành tính (\texttt{DROP}/\texttt{LOG}), nên một dị biệt đã gắn cờ không bao giờ làm nhiễm thống kê nó đang được đo---trực tiếp chống đầu độc ``ếch luộc'' tốc độ chậm. Thứ hai, thay vì khởi động lạnh từ warmup có thể bị nhiễm, engine có thể được seed từ một \emph{golden baseline}: trạng thái Welford ($n$, trung bình, tổng-bình-phương) được tính trước offline trên bản ghi benign đã kiểm định, nên dị biệt được đo trên thống kê sạch-đã-biết ngay từ gói đầu, sau đó cùng cập nhật có điều kiện tiếp tục.

\section{Cổng ML: Giảm tải Tầng 1 bằng Học máy}
Vì đẩy mọi ca leo thang Tầng 1 lên một LLM nhiều giây là lãng phí khi nhiều ca giải được bằng thống kê, một cổng \texttt{LightGBM} được đặt giữa bộ luật và tầng nhận thức. Huấn luyện offline trên $\approx\!1$M luồng có nhãn (Test-F1 held-out $=0.9635$, xuất dưới dạng \texttt{tier\_2\_model.pkl} chứa scaler và các đặc trưng), nó chấm xác suất tấn công $C$ của mỗi ca leo thang và áp dụng chính sách độ tin cậy bốn dải thống nhất cũng chi phối phán quyết LLM: $C\!\ge\!0.85\rightarrow\texttt{BLOCK\_IP}$ (kèm sinh luật động), $0.65\!\le\!C\!<\!0.85\rightarrow$ leo thang lên LLM, $0.40\!\le\!C\!<\!0.65\rightarrow\texttt{ALERT}$ (tự chặn nếu tái phạm), và $C\!<\!0.40\rightarrow\texttt{DROP}$. Để chống né tránh (anti-evasion), hệ thống kẹp Z-score tại $\pm8\sigma$ và điền \texttt{NaN} bằng giá trị trung bình; nếu hơn $30\%$ đặc trưng vượt $6\sigma$, mô hình tự động abstain và đẩy thẳng lên LLM. Suy luận tốn $\approx\!0.35$\,ms, nên cổng tự quyết phần lớn ca chắc chắn ở tốc độ đường truyền và chỉ chuyển tiếp dải bất định; thực nghiệm nó giảm tải $83.8\%$ ca leo thang ở precision $98.82\%$ (Chương~\ref{ch:experiments_evaluation}), biến LLM thành tài nguyên khan hiếm được dành riêng chứ không phải chi phí bắt buộc mỗi sự kiện. Chia sẻ cùng module \texttt{decision\_policy} giữa tầng ML và LLM giữ ánh xạ độ-tin-cậy-tới-hành-động nhất quán ở cả hai tầng.

\section{Tầng 2: Tác tử Nhận thức Có trạng thái LangGraph}
Các ca leo thang bất định còn lại vào engine Tầng 2, một Đồ thị Có hướng Không chu trình được LangGraph~\cite{langgraph2024} biên dịch và hoạt động như một máy trạng thái hữu hạn trên đối tượng \texttt{SentinelState} trung tâm. Đồ thị nối một nút rào chắn (nén + đóng gói phân định), một nút RAG-context, một nút LLM-triage, và một nút ánh xạ ATT\&CK có điều kiện, kết thúc ở bộ thực thi hành động hoặc bộ xử lý con-người-trong-vòng-lặp. Nút RAG truy xuất nền \emph{vô điều kiện} trước triage. Nút triage đánh giá lô bằng mô hình Gemma-2 lượng tử hóa cục bộ~\cite{gemma_2_2024} và phát một phán quyết JSON có cấu trúc; mọi phán quyết hành động (\texttt{BLOCK\_IP}, \texttt{ALERT}, \texttt{AWAIT\_HITL}) được định tuyến qua bộ ánh xạ, nơi neo vào kỹ thuật triage đã nối nền và cấu trúc hóa thành bản ghi MITRE kiểm chứng schema, trong khi phán quyết \texttt{LOG} lành tính kết thúc. Định tuyến là tất định và không chu trình (Bảng~\ref{tab:agent_fsm_transitions}).

\begin{table}[!htbp]
    \centering
    \caption{Ma trận chuyển trạng thái của FSM tác tử nhận thức Tầng 2. (Nguồn: Tác giả)}
    \label{tab:agent_fsm_transitions}
    \begin{tabular}{|>{\raggedright\arraybackslash}p{0.20\textwidth}|>{\raggedright\arraybackslash}p{0.27\textwidth}|>{\raggedright\arraybackslash}p{0.18\textwidth}|>{\raggedright\arraybackslash}p{0.16\textwidth}|}
        \hline
        \textbf{Trạng thái} & \textbf{Đầu vào / Điều kiện} & \textbf{Kế tiếp} & \textbf{Hành động} \\
        \hline
        \texttt{GUARDRAILS} & Nhận lô & \texttt{RAG\_CONTEXT} & Nén \& bọc nonce \\ \hline
        \texttt{RAG\_CONTEXT} & Truy xuất ngữ cảnh & \texttt{LLM\_TRIAGE} & Thêm MITRE/NIST \\ \hline
        \texttt{LLM\_TRIAGE} & action hành động & \texttt{ATTACK\_MAPPER} & Cấu trúc ATT\&CK \\ \hline
        \texttt{LLM\_TRIAGE} & action $\in$ \{\texttt{LOG},\texttt{DROP}\} & \texttt{END} & Đóng chu kỳ \\ \hline
        \texttt{ATTACK\_MAPPER} & \texttt{BLOCK\_IP}/\texttt{ALERT} & \texttt{ACTION\_EXEC} & Chặn / cảnh báo \\ \hline
        \texttt{ATTACK\_MAPPER} & \texttt{AWAIT\_HITL} & \texttt{HITL} & Đưa analyst \\ \hline
        \texttt{ACTION\_EXEC}/\texttt{HITL} & Hoàn tất & \texttt{END} & Ghi bản ghi HMAC \\ \hline
    \end{tabular}
\end{table}

Nối nền triage trong tài liệu MITRE/NIST truy xuất ràng buộc LLM và giảm ảo giác. Một \texttt{DecisionValidator} tất định hạ mọi \texttt{BLOCK\_IP} dưới độ tin cậy $0.5$ về \texttt{AWAIT\_HITL} và ghi đè LLM thành \texttt{AWAIT\_HITL} bất cứ khi nào nó mâu thuẫn tín hiệu tấn công Tầng 1 (rào chắn đồng thuận chống thao túng ngữ nghĩa). Thêm vào đó, một \textbf{T1571 Guardrail} ép các tấn công dựa trên cổng chung chung (MITRE T1571) về \texttt{AWAIT\_HITL} để tránh chặn nhầm. Để tối ưu luồng, nếu \texttt{attack\_mapper} không khớp hoặc không rõ ràng, SENTINEL bỏ qua lần gọi LLM thứ 2 và mặc định thẳng về \texttt{AWAIT\_HITL}. \texttt{LLMClient} kết nối dịch vụ llama.cpp cục bộ chạy Gemma-2 với \texttt{temperature=0.1}, seed cố định, cửa sổ 8{,}192 token, và giải mã Pydantic (\texttt{LLMDecision}) JSON ràng buộc schema. Nếu chuỗi JSON do LLM sinh ra bị lỗi định dạng, client sử dụng cơ chế dự phòng dựa trên regex (\texttt{\_salvage\_fields}) để khôi phục dữ liệu, và mặc định trả về \texttt{AWAIT\_HITL} nếu việc khôi phục thất bại, qua đó đảm bảo khả năng chịu lỗi (fault tolerance) vững chắc. Cơ chế KV-Caching được tận dụng triệt để: system prompt giữ ở dạng tĩnh trong khi feedback biến động được chuyển vào user message, cho phép llama.cpp cache prompt hiệu quả và giảm TTFT.

\section{Cơ chế Dual-RAG và Bộ nhớ Ngữ nghĩa}
Dual-RAG nối nền quyết định bằng cách truy xuất từ hai kho: MITRE ATT\&CK~\cite{strom2018mitre} (tactic/technique) và NIST SP~800-61r2~\cite{nist80061} (playbook ứng phó). Thay vì truy xuất dày hay thưa riêng lẻ~\cite{rag2020}, nó hợp nhất cả hai (Bảng~\ref{tab:search_comparison}).

\begin{table}[!htbp]
    \centering
    \caption{Các mô hình truy xuất trong chiến lược tìm kiếm lai. (Nguồn: Tác giả)}
    \label{tab:search_comparison}
    \begin{tabular}{|>{\raggedright\arraybackslash}p{0.20\textwidth}|>{\raggedright\arraybackslash}p{0.31\textwidth}|>{\raggedright\arraybackslash}p{0.31\textwidth}|}
        \hline
        \textbf{Chiều} & \textbf{Dày đặc (FAISS)} & \textbf{Thưa (BM25)} \\
        \hline
        Mô hình & all-MiniLM-L6-v2, 384 chiều~\cite{reimers2019sentence} & Trọng số hạng BM25 \\ \hline
        Logic & Tương đồng cosine~\cite{faiss2019} & TF-IDF~\cite{bm25_2009} \\ \hline
        Điểm mạnh & Đồng nghĩa, khái niệm, ngữ nghĩa & ID chính xác (CVE, IP, cổng) \\ \hline
        Hạn chế & Yếu với hash kỹ thuật chính xác & Mù đồng nghĩa/ngữ cảnh \\ \hline
    \end{tabular}
\end{table}

Điểm dày (FAISS) và thưa (BM25) được hợp nhất bằng Reciprocal Rank Fusion~\cite{rrf2009},
\begin{equation}
    R(d) = \sum_{m \in M} \frac{1}{k + r_m(d)},
\end{equation}
với $r_m(d)$ là hạng của $d$ trong run $m$ và $k=60$; ba đoạn hợp nhất hàng đầu nối nền cho nút triage. Một Semantic Cache Tầng 1.75 đứng trước truy xuất, khóa theo băm SHA-256 của mẫu log đã khai thác nên các mẫu tấn công lặp (vd: DDoS, Brute-force) tái dùng kết quả đã tính và bỏ qua LLM hoàn toàn, giảm độ trễ quyết định đầu-cuối xuống $\approx\!1$\,ms. Nó là LRU trong bộ nhớ (\texttt{OrderedDict}) giới hạn 500 mục với TTL 1800 giây.

\section{Bộ nhớ Đe dọa và Liên kết Chiến dịch APT}
Kẻ tấn công ẩn mình chia chuỗi tấn công thành sự kiện độ-nghiêm-trọng-thấp rải qua cửa sổ dài, vượt bộ lọc phi trạng thái~\cite{apt_stealth2021}. SENTINEL thêm Bộ nhớ Đe dọa bền vững (SQLite cục bộ) theo dõi cảnh báo lịch sử, uy tín thực thể, và chuỗi thời gian (Bảng~\ref{tab:threat_memory_schema}). CSDL SQLite được cấu hình \texttt{synchronous=NORMAL} và có đánh chỉ mục, nhưng tuyệt đối tránh chế độ \texttt{WAL} (\texttt{journal\_mode=WAL}) để ngăn lỗi crash phân quyền cross-UID (liên quan đến tệp \texttt{-wal} và \texttt{-shm}) trong môi trường Docker.

\begin{table}[!htbp]
    \centering
    \caption{Lược đồ cơ sở dữ liệu Bộ nhớ Đe dọa. (Nguồn: Tác giả)}
    \label{tab:threat_memory_schema}
    \begin{tabular}{|>{\raggedright\arraybackslash}p{0.24\textwidth}|>{\raggedright\arraybackslash}p{0.34\textwidth}|>{\raggedright\arraybackslash}p{0.28\textwidth}|}
        \hline
        \textbf{Bảng} & \textbf{Trường chính} & \textbf{Mục đích} \\
        \hline
        \texttt{ip\_reputation} & \texttt{ip} (PK), \texttt{reputation\_score}, \texttt{total\_incidents}, \texttt{last\_seen}, \texttt{last\_mitre\_technique} & Uy tín IP dài hạn (0--100), suy giảm theo thời gian. \\ \hline
        \texttt{known\_entities} & \texttt{id} (PK), \texttt{entity\_type}, \texttt{entity\_value}, \texttt{is\_active} & Seed sẵn công cụ nội bộ đã kiểm định để ép báo giả. \\ \hline
        \texttt{threat\_events} & \texttt{id} (PK), \texttt{src\_ip}, \texttt{apt\_phase}, \texttt{apt\_day}, \texttt{label}, \texttt{timestamp} & Ghi sự kiện leo thang cho liên kết APT đa ngày. \\ \hline
        \texttt{apt\_indicators} & \texttt{id} (PK), \texttt{indicator\_value}, \texttt{confidence}, \texttt{mitre\_chain} & Chỉ báo APT xác nhận cho giám sát chiến dịch. \\ \hline
    \end{tabular}
\end{table}

Uy tín một IP tăng theo độ nghiêm trọng cảnh báo, tối đa 100; một hàm suy giảm định kỳ giảm điểm của IP im lặng,
\begin{equation}
    S_{new} = S_{old} \times (1 - \lambda)^{t},
\end{equation}
với tỷ lệ suy giảm $\lambda$ và thời gian im lặng $t$. Một routine APT quét \texttt{threat\_events}: nếu một IP nguồn kích hoạt dị biệt qua $\ge\!2$ ngày riêng biệt (tùy chọn trải nhiều pha kill-chain), phân loại được nâng lên chiến dịch đa ngày nổi lên. Các khoảng và gia số này được tính bằng truy vấn upsert/tổng hợp phía Python trên SQLite, tránh giữ trạng thái dài hạn trong bộ nhớ làm việc của tác tử. Tầng 1 cũng đọc uy tín này, nên một IP tiền sử thù địch (vd: $\texttt{reputation}=100$) bị chặn ngay lập tức (block-on-sight) không cần gọi LLM thêm---khép vòng phản hồi. Ngược lại, các IP hạ tầng nội bộ (vd: \texttt{10.0.0.99}) được hardcode trong whitelist và miễn nhiễm tuyệt đối với việc chặn.

\section{Rào chắn Bảo mật AI và Lớp Đóng gói}
Vì LLM dễ bị chiếm quyền chỉ dẫn~\cite{promptinjection2023}, một rào chắn nhiều lớp cách ly engine suy luận. Phòng thủ lõi là \emph{Đóng gói Dữ liệu Phân định}: mỗi session sinh một nonce mật mã (vd \texttt{<<<DATA\_HASH...>>>}), bọc payload không tin cậy để mọi token lệnh bên trong biên được coi là dữ liệu thụ động. Hơn nữa, việc phát hiện prompt injection được tối ưu hóa bằng so khớp chuỗi con $O(1)$ thay vì thực thi regex đắt đỏ. Chống token-DoS và tràn ngữ cảnh, một bộ khai thác mẫu Drain~\cite{he2017drain} gộp log tương tự và loại tham số ngẫu nhiên trong ngân sách 4{,}000 token (60\% mẫu, 40\% ngoại lai entropy cao). Một bộ lọc làm sạch đầu ra ép tuân thủ schema JSON và loại vector markdown-image, SVG, hex-bypass~\cite{owasp_llm2025} trước khi đầu ra tới cơ sở dữ liệu hay dashboard.

\section{Tính toàn vẹn Dữ liệu và Bảo vệ HMAC}
Để bảo vệ cơ sở dữ liệu log khỏi giả mạo, SENTINEL móc xích từng bản ghi bằng mật mã~\cite{hmac1997}: với bản ghi $i$,
\begin{equation}
    H_i = \text{HMAC-SHA256}(D_i \parallel H_{i-1}, K),
\end{equation}
với $D_i$ là payload hiện tại, $H_{i-1}$ là băm trước, và $K$ là khóa bí mật. Sửa một nhãn hay xóa một mục làm đứt chuỗi, được dashboard phát hiện ngay. Kho vector RAG cũng được một bộ quản lý checksum SHA-256 bảo vệ, xác minh tệp trước khi nạp, chặn đầu độc ngữ nghĩa. Một \texttt{ActionValidator} ép danh sách cho phép các hành động hợp lệ và từ chối ký tự shell trong trường payload, ngăn RCE khi điều phối tường lửa.

\section{Dashboard, Triển khai và Khả năng Mở rộng}
Dashboard SOC Streamlit đọc phản ứng cơ sở dữ liệu và tệp backend, tích hợp giao diện hàng đợi thời gian thực (băng thông Tầng 1, phán quyết Tầng 2), bảng phân tích Bộ nhớ Đe dọa (lan truyền APT theo dòng thời gian IP), chỉ báo toàn vẹn audit HMAC (xanh = chuỗi hợp lệ, đỏ = giả mạo), và giao diện HITL để duyệt hoặc bác các chặn đề xuất. Triển khai được ảo hóa bằng \texttt{Dockerfile} đa giai đoạn (llama.cpp biên dịch với CUDA cho suy luận GPU) điều phối bởi \texttt{docker-compose.yml}, chia sẻ thư mục log/CSDL qua volume bền vững nên \texttt{threat\_memory.db} và cơ sở dữ liệu audit sống sót qua vòng đời container. Quyền riêng tư dữ liệu được tuân thủ nghiêm ngặt: mật khẩu trong \texttt{REDIS\_URL} bị che (redact) trước khi log, và nhãn \texttt{apt\_mitre\_ttp} bị gỡ (strip) khỏi log trước khi đưa vào LLM (leak guard). Cuối cùng, để xử lý luồng sự kiện lớn, hệ thống quản lý tải bằng độ trễ (\texttt{lag}) của Redis consumer group thay vì độ dài hàng đợi (\texttt{xlen}), cung cấp đồng bộ backpressure bền bỉ.

Thiết kế hai tầng trả lời phản biện băng thông bằng cách coi hai tầng là hai chế độ mở rộng khác nhau. Tầng 1 là bài toán băng thông, không phải bức tường độ trễ: mỗi cập nhật Welford và khớp regex là $O(1)$ và trạng thái per-IP độc lập, nên lớp này song song hóa dễ dàng và phân mảnh theo băm IP nguồn; mỗi quyết định chạy dưới một mili-giây (Cổng ML thêm $\approx\!0.35$\,ms), và một triển khai biên dịch hoặc stream-processing nâng trần mỗi lõi nhiều bậc độ lớn. Tầng nhận thức không bao giờ thấy số leo thang thô: nén Drain, Semantic Cache, và bỏ ở Tầng 1 theo uy tín thu gọn nó về số mẫu dị biệt \emph{thực sự mới}, phán xử theo lô ($\approx\!5.7$\,s mỗi lô, một lệnh suy luận cho hàng chục ca leo thang nén). Quan trọng, tầng nhận thức nằm \emph{ngoài} đường đồng bộ: Tầng 1 đã chặn hoặc bỏ tấn công lộ rõ ở tốc độ đường truyền, nên tồn đọng nhận thức làm chậm làm-giàu chứ không phải bảo vệ. Dưới một trận lụt đối kháng toàn-mới, chính đợt tăng là một ngoại lai volumetric Tầng 1 kích hoạt giới hạn tốc độ thô, và một hàng đợi nhận thức bão hòa lùi về phán quyết Tầng 1 tất định---bảo vệ được giữ, chỉ làm-giàu suy biến. Trung thực mà nói, bản mẫu hiện tại triển khai tỷ lệ giảm tải và caching nhưng chưa có bộ điều khiển backpressure tường minh; hình thức hóa chính sách suy biến duyên dáng đó, mở rộng GPU ngang, và một kho golden-baseline tái kiểm định định kỳ là con đường tới sản xuất, được xem lại ở Chương~\ref{ch:conclusion}.

\section{Tóm tắt Chương}
Chương này đã thiết lập thiết kế SENTINEL: lọc phi trạng thái tốc độ đường truyền (Tầng 1), một Cổng ML học máy giảm tải phần lớn ca leo thang chắc chắn, và suy luận nhận thức có trạng thái (Tầng 2), tất cả bọc trong rào chắn mật mã và đặc-thù-AI cung cấp phòng thủ chiều sâu---giảm chứ không loại bỏ phơi nhiễm đối kháng, với độ bền trước kẻ tấn công thích ứng là thuộc tính giới hạn xem lại ở Chương~\ref{ch:conclusion}. Dashboard Streamlit nối tự động hóa với giám sát con người qua hàng đợi HITL, và cấu hình Docker bảo đảm triển khai tái lập được. Chương~\ref{ch:experiments_evaluation} trình bày đánh giá thực nghiệm kiến trúc này.
